\section{Binary Search: a Lemma over a Whole Cycle}\label{sec:bisect}

Binary search is the standard example of a loop whose variables are only meaningful together: $lo$ and $hi$
narrow an interval through a probe computed from both. This section shows what it takes to prove it by lookup:
a lemma that names the loop as a whole, two facts about the named loop, and the vocabulary of sorted arrays
that turns the facts into the postcondition.

\subsection{The programs}
\begin{lstlisting}
method search(a: array<int, int>, n: int, x: int) returns (i: int)
  requires 0 <= n && sorted(a, 0, n)
  ensures 0 <= i ==> i < n && a[i] == x
  ensures i < 0 ==> forall k: int :: 0 <= k < n ==> a[k] != x
{
  var lo: int; var hi: int; var mid: int;
  lo := 0; hi := n; i := -1;
  while lo < hi {
    mid := (lo + hi) / 2;
    if a[mid] < x { lo := mid + 1; }
    else if x < a[mid] { hi := mid; }
    else { i := mid; break; }
  }
}
\end{lstlisting}
The second program, the lower bound, has no third branch: \texttt{else \{ hi := mid; \}}, and returns $lo$ with the
postconditions $0 \le lo \le n$, $a[k] < x$ for $k < lo$ and $x \le a[k]$ for $lo \le k < n$. Both take the
array sorted on $[0, n)$ and nothing else.

\subsection{The system}
The two variables form one trace of two accumulators. With $M(lo, hi) = (hi + lo)/2$, $G(m) = a[m] < x$ and
$H(m) = x < a[m]$,
\begin{align}
  S &= \mathsf{trace}(0, n,\ \mathsf{for}\ \bl{K}\ \mathsf{do}\ lo, hi, k \mapsto
     \upd{G(M)}{M + 1}\ \ \upd{H(M)}{lo}, \nonumber\\
  &\qquad\qquad \upd{G(M)}{hi}\ \ \upd{H(M)}{M}) \label{eq:bs1}\\
  &\al{\bl{K}, \bl{S}.0, \bl{S}.1}{k, lo, hi \mapsto lo < hi} \label{eq:bs2a}\\
  &\al{\bl{K}, \bl{S}.0, \bl{S}.1}{k, lo, hi \mapsto G(M) \vee H(M)} \label{eq:bs2}
\end{align}
with $lo = \cur(S).0$, $hi = \cur(S).1$. The second constraint \eqref{eq:bs2} is there because of the \texttt{break}: an iteration
that finds $x$ leaves the loop, so every iteration that completed took one of the two other branches. After the
loop there are two paths: the normal exit with $hi \le lo$ and $i = -1$, and the \texttt{break} with $lo < hi$,
$a[M(lo, hi)] = x$ and $i = M(lo, hi)$.

\subsection{What the proof needs}
\paragraph{A name for the loop.} The database defines the bisection of an interval with an arbitrary probe and
two guards, and folds it where its body occurs:
\begin{lstlisting}
function {:teq} {:fold} bisection(l0: int, h0: int, ks: List<int>, M: (int, int) -> int,
                                  G: int -> bool, H: int -> bool): Trace<(int, int)>
{
  trace(l0, h0, for ks do lo, hi, k ::
        case G(M(lo, hi)) && lo < hi => M(lo, hi) + 1
        case H(M(lo, hi)) && !G(M(lo, hi)) && lo < hi => lo
        case !(lo < hi) || !(G(M(lo, hi)) || H(M(lo, hi))) => lo,
        case G(M(lo, hi)) && lo < hi => hi
        case H(M(lo, hi)) && !G(M(lo, hi)) && lo < hi => M(lo, hi)
        case !(lo < hi) || !(G(M(lo, hi)) || H(M(lo, hi))) => hi)
}
\end{lstlisting}
The fold matches the trace \eqref{eq:bs1} as a whole (Section~\ref{sec:engine:cycles}). The probe $M$ is bound from
the result $M(lo, hi) + 1$ by abstraction, $G$ and $H$ from the guards; the conjunct $lo < hi$ of the definition's
guards follows from the first constraint \eqref{eq:bs2a}, which holds inside the update, and $\neg G(M)$ in the
second case follows from the program's guard $x < a[m]$. The third case covers no case of the program and is
ignored (Section~\ref{sec:engine:cases}). The result is $S = \mathsf{bisection}(0, n, \bl{K}, M, G, H)$.

\paragraph{Two facts about the named loop.} If $G$ holds on a prefix of $[l_0, h_0)$, the lower end stays within $l_0$
and the \emph{cut} of $G$, the first index where it fails; if $H$ holds on a suffix, the upper end stays within the
cut of $\neg H$ and $h_0$:
\begin{lstlisting}
lemma {:fw} {:teq} bisection_lo(l0: int, h0: int, ks: List<int>, M: (int, int) -> int,
                                 G: int -> bool, H: int -> bool)
  requires l0 <= h0
  requires forall l: int, h: int :: l < h ==> l <= M(l, h) && M(l, h) < h
  requires forall p: int, q: int :: l0 <= p && p <= q && q < h0 && G(q) ==> G(p)
  ensures l0 <= current(bisection(l0, h0, ks, M, G, H)).0 <= cut(G, l0, h0)
\end{lstlisting}
and \lemma{bisection\_hi} symmetrically, $\mathsf{cut}((j) \Rightarrow \neg H(j), l_0, h_0) \le \cur(\ldots).1 \le h_0$
under the condition that $H$ holds on a suffix. They are forward facts, triggered by $\cur(S).0$ and $\cur(S).1$ in the
goal. Their conditions are decided after the instantiation: the probe lies strictly inside the interval by the
bounds of division (Section~\ref{sec:engine:arith}), and the guards hold on a prefix, respectively a suffix, since
the rules below turn the comparisons $a[q] < x$ for indices in range into comparisons with a partition point,
$q < \mathsf{lower}(a, n, x)$, after which the condition is linear.

\paragraph{The vocabulary of sorted arrays.} The cuts become partition points:
$\mathsf{lower}(a, n, x)$, below which the elements are less than $x$, and $\mathsf{upper}(a, n, x)$, from which
they are greater: \lemma{cut\_lower} rewrites the cut of $a[j] < x$, and the cut of $\neg(x < a[j])$, i.e.\ of
$a[j] < x + 1$, is $\mathsf{lower}(a, n, x + 1) = \mathsf{upper}(a, n, x)$ (\lemma{lower\_succ}). For an index in range,
the comparisons with $x$ are comparisons with the partition points (\lemma{lower}, \lemma{upper},
$a[m] = x \Leftrightarrow \mathsf{lower} \le m < \mathsf{upper}$ by \lemma{equal}), and the partition points lie within
$0 \le \mathsf{lower} \le \mathsf{upper} \le n$ (the forward fact \lemma{bounds}).

\paragraph{The postconditions.} At the normal exit, $\mathsf{upper} \le hi \le lo \le \mathsf{lower}$, so the range
$[\mathsf{lower}, \mathsf{upper})$ of the elements equal to $x$ is empty and no $a[k]$ equals $x$. At the
\texttt{break}, $a[M] = x$ and $M < hi \le n$. For the lower bound, $H$ is $x \le a[m]$, the negation of $G$; the cut
of $\neg H$ is the cut of $G$, so $lo \le \mathsf{lower} \le hi$, and at the exit $hi \le lo$ gives
$lo = \mathsf{lower}(a, n, x)$, from which the two quantified postconditions follow by \lemma{lower} and
\lemma{lower\_le}. Both programs close in under a minute.

\subsection{Why the definition has three cases}
The program's update has two cases. A definition with only those two, guarded by $lo < hi$, would be
\emph{unspecified} in an iteration where $lo < hi$ fails or neither guard holds, and \lemma{bisection\_lo}, stated for
any number of steps, would bound that unspecified value. Take $l_0 = h_0$, one step, and $G$, $H$ false: all
conditions of \lemma{bisection\_lo} hold, the cut of $G$ on the empty interval is $l_0$, and the lemma says that the
first component of $\mathsf{unspecified}$ equals $l_0$, for every $l_0$, which derives $\mathsf{false}$. The facts
are true of the loops that the programs run, whose iterations all satisfy the constraints, but not of every trace
that the definition denotes.

The third case removes the gap: where the program's step is undefined, the definition keeps the state, which keeps
the bounds. The bounds now hold for every run of the definition, and they are proved by induction over the steps,
carrying $l_0 \le lo \le \mathsf{cut}(G)$ and $hi \le h_0$ (for $hi$: $l_0 \le lo$ and $\mathsf{cut}(\neg H) \le hi \le h_0$).
The second case excludes $G$, so that the guards are disjoint: where $G$ and $H$ both hold the first case is taken,
which the bound also allows, since $G(m)$ and $G$ holding on a prefix give $m < \mathsf{cut}(G)$. The fold still
matches the program because the extra case is excluded by the context of the update.

\paragraph{Why $H$ cannot be replaced by $\neg G$.} In the context of the update, the program's guards are disjoint
and complete, so $H(m)$ and $\neg G(m)$ agree at every probe that the loop visits. A definition with $G$ alone,
moving $hi$ down wherever $G$ fails, matches the programs as well, but its fact about $hi$ can only be
$\mathsf{cut}(G) \le hi$: a run of that definition may move $hi$ onto a position that holds $x$, which the search
leaves by the \texttt{break}. The search needs $\mathsf{upper} \le hi$, i.e.\ the cut of $\neg H$, a property of $H$
on the whole interval rather than at the probes. The definition must therefore encode the exit itself: a probe at
which neither guard holds keeps the state, and $hi$ moves only to positions where $H$ holds.
